\chapter{力学}

位形空间$X$通常是一个$n$-维光滑流形，
其上的函数集构成光滑函数代数$C(X)$。
切丛$TX$构成物理上的\textbf{速度相空间}(velocity phase space)，
在局部坐标系表示的广义坐标$x^i$下，
有广义速度$\dot q^i = \text d q^i / \text d t$，
二者构成$TX$的局部坐标。
考虑函数$L(q^i,\dot q^i,t)$，
它在时间上的积分称为\textbf{作用量}(action)：

$$
S = \int_0^T L \text d t
$$

\textbf{Hamilton变分原理}(variational principle of Hamilton)要求作用量最小，
这样的极值问题涉及到对泛函求变分使得：

$$
\delta S = \delta \int_0^T L \text d t = 0
$$

需要以独立的方式去扰动$C(X)$的点。
函数求导扰动的是用实数无穷小量$\varepsilon$扰动$x$：

$$
\frac{\text d f}{\text d x} \bigg|_{x = x_0}
= \lim_{\varepsilon \to 0}
\frac{f(x_0 + \varepsilon) - f(x_0)}{\varepsilon}
$$

而泛函更复杂，
如果上述极限的分母仍是实数无穷小量$\varepsilon$，
那么$\varepsilon$扰动$f$需要引入一个函数$g$来产生
$f(x) + \varepsilon g(x)$。
普通的$g$会同时扰动$f$在许多$x$处的函数值，
无法独立出点$x=x_0$对函数变换的贡献。
为了在$x = x_0$点附近扰动，
狄拉克$\delta$函数就很适合：

$$
\delta(x) = \begin{cases}
+\infty, & x=0,\\
0, & x\ne 0,
\end{cases}
\qquad \text{s.t.} \quad
\int_{-\infty}^{\infty}\delta(x)\,dx = 1.
$$

于是$f(x) + \varepsilon \delta$只在$x = x_0$处扰动了$f$。
应用狄拉克$\delta$函数构造\textbf{泛函导数}(functional derivative)：

$$
\frac{\delta F}{\delta f}
= \lim_{\varepsilon\to 0}
\frac{F[f(x)+\varepsilon\delta(x^\prime - x)]-F[f(x)]}{\varepsilon}
$$

现在需要能让作用量最小的Lagrangian $L$，
满足：

$$
\frac{\delta L}{\delta x} = 0
$$

总机械能$E = T + V$守恒的情况下，
动能和势能此消彼长。
考虑到量子力学只能测量平均值，
在时间区间$[0,\tau]$上通过积分得到

$$
\overline T = \frac{1}{\tau}\int_0^\tau \frac{m \dot x^2 \text d t}{2}, \quad 
\overline V = \frac{1}{\tau}\int_0^\tau V[x(t)] \text d t
$$

求二者的变分\cite{Lancaster2014}：

$$
\frac{\delta \overline T[x]}{\delta x(t)} = - \frac{m \ddot x}{\tau}, \quad
\frac{\delta \overline V[x]}{\delta x(t)} = \frac{V^\prime[x(t)]}{\tau}
$$

于是

$$
\frac{\delta (\overline T[x] - \overline V[x])}{\delta x(t)} = 0
$$

这样得到了我们需要的Lagrangian:

$$
L = T - V
$$

满足\textbf{Hamilton最小作用量原理}(Hamilton's principle of least action)。
视坐标和速度为独立变量:

$$
\frac{\delta S}{\delta x(t)} 
= \frac{\delta L}{\delta x(t)} - \frac{\text d}{\text d t} \frac{\delta L}{\delta \dot x(t)} 
= 0
$$

得到Euler-Lagrange方程。


\cite{MarsdenRatiu1994}


